% Implementation
\section{Implementation}
\label{sec:implementation}

This section describes the practical realization of SPOQ, including the file-system
representation of epics and tasks, the journal-based tracking system for agent
work sessions, and the skill framework that encapsulates reusable agent capabilities.
See Appendix~\ref{subsec:implementation-detail} for complete schema examples,
journal entry format, and skill metadata.

\subsection{Task Representation: The spoq/epics/ Directory}
\label{subsec:task-representation}

SPOQ represents epics and tasks as structured files in a designated directory,
enabling version control, human review, and programmatic manipulation.

\paragraph{Epic Lifecycle.}
Epics follow a two-stage directory lifecycle that reflects their progression through
the SPOQ pipeline. During planning, new epics are created under
\texttt{spoq/epics/active/}, where they remain throughout validation and execution
phases. Once every constituent task passes the agent-validation gate, the orchestrator
relocates the entire epic directory to \texttt{spoq/epics/complete/}. This
filesystem-level transition serves as an unambiguous completion signal: version
control history preserves the move as a single commit, providing an auditable
record of when each epic achieved full validation. Separating in-flight work from
verified deliverables also prevents agents from inadvertently modifying artifacts
that have already satisfied quality thresholds.

\newpage
\paragraph{Roadmap Coordination.}
The \texttt{ROADMAP.md} file at the root of \texttt{spoq/epics/} acts as a
centralized registry of all epics and their current disposition. During planning,
the orchestrator appends a new entry containing the epic identifier, a brief
objective summary, and an initial status of \emph{planned}. As execution proceeds,
status fields are updated to reflect transitions through \emph{in-progress},
\emph{validation}, and \emph{done}. Because the roadmap records inter-epic
dependencies alongside status, it enables the orchestrator to determine which
downstream epics become unblocked when a predecessor completes.

\paragraph{Task YAML Schema.}
Each task file follows a standardized YAML schema with three categories of fields:

\begin{definition}[Task Specification]
\label{def:task-spec}
A \textit{task specification} is a YAML document with the following structure:
\end{definition}

\textbf{Identity fields} (\texttt{id}, \texttt{title}, \texttt{epic}) establish
task context by assigning a unique identifier within the epic, a human-readable
title, and a back-reference to the parent epic.
%
\textbf{Execution control fields} govern scheduling through status tracking
(\texttt{pending}/\texttt{in\_progress}/\texttt{completed}), priority levels,
wave assignment via the \texttt{phase} field, PERT three-point effort estimates,
dependency lists referencing prerequisite task IDs, and required domain skills.
%
\textbf{Deliverable fields} define verification criteria: \texttt{files\_to\_touch}
enumerates all files the agent may modify, \texttt{outputs} lists tangible
deliverables, and \texttt{acceptance\_criteria} provides a checklist the agent
must satisfy before marking the task complete.
%
A free-form \textbf{description} field supplies structured implementation
guidance using Markdown, typically containing an objective statement,
step-by-step instructions with code snippets, and verification commands.
See Appendix~\ref{subsec:implementation-detail} for annotated examples of each
field category.

\paragraph{Phase Assignment.}
The \texttt{phase} field encodes wave assignment from topological analysis:
\begin{itemize}
    \item Phase 0: Tasks with no dependencies (Wave 0)
    \item Phase $n$: Tasks depending only on tasks in phases $< n$
\end{itemize}

Tasks within the same phase execute concurrently. The critical path determines
the minimum number of waves required.

\subsection{Journal Tracking System}
\label{subsec:journal}

SPOQ employs a journal-based tracking system that records agent work sessions
with structured metadata, enabling explainability, meta-orchestration, and
training data generation.

\begin{definition}[Journal Entry]
\label{def:journal-entry}
A \textit{journal entry} consists of YAML frontmatter followed by Markdown
sections documenting a completed work session.
\end{definition}

The frontmatter captures machine-readable metadata: agent identity, ISO 8601
timestamps, a calibrated confidence score (0.0--1.0), session type classification,
the list of modified files, and task completion counts. The Markdown body follows
a standardized layout with sections for summary, work completed, changes made,
issues encountered, testing results, and next steps. See
Appendix~\ref{subsec:implementation-detail} for complete format examples and
tooling support.

\paragraph{Confidence Scoring.}
Each journal entry includes a calibrated confidence score (0.0--1.0) reflecting
the agent's self-assessment of work quality. Scores above 0.85 indicate
well-tested, production-ready output; scores between 0.65--0.84 signal
functional work requiring additional validation; scores below 0.65 flag
experimental results with known gaps. See Table~\ref{tab:confidence} in the
Appendix for the full interpretation scale.

\paragraph{Value Proposition.}
The journal system provides four categories of value:

\begin{enumerate}
    \item \textbf{Explainability (XAI):} Every work session is documented with
    rationale, enabling audit trails and decision archaeology.

    \item \textbf{Multi-Agent Coordination:} Timestamps and task progress enable
    parallel agents to avoid conflicts and build on each other's work.

    \item \textbf{Knowledge Graph Construction:} Journal entries create nodes
    (sessions, files, tasks) and edges (MODIFIED, COMPLETED) for analysis.

    \item \textbf{Training Data:} Combined with git commits, journal entries
    provide rich examples for fine-tuning development-focused LLMs.
\end{enumerate}

\subsection{Skill Framework}
\label{subsec:skills}

SPOQ encapsulates reusable agent capabilities as \textit{skills}, self-contained
modules that provide domain knowledge, workflows, and tooling for specific
task categories.

\begin{definition}[Skill]
\label{def:skill}
A \textit{skill} is a directory containing:
\begin{itemize}
    \item \texttt{SKILL.md}: Instructions and metadata (required)
    \item \texttt{scripts/}: Executable utilities (optional)
    \item \texttt{references/}: Documentation loaded on demand (optional)
    \item \texttt{assets/}: Templates and configurations (optional)
\end{itemize}
\end{definition}

\paragraph{Core SPOQ Skills.}
The framework includes six skills aligned with the four-stage pipeline
(see Table~\ref{tab:skills} in the Appendix for the full inventory).

\paragraph{Skill Invocation and Anatomy.}
Skills are invoked via slash commands (e.g., \texttt{/epic-planning},
\texttt{/agent-execution}) that expand into full prompts with context from the
referenced epic directory. Each skill's \texttt{SKILL.md} contains YAML
frontmatter (name, description) followed by structured documentation covering
activation criteria, core patterns, quality standards, and integration points
with other skills. See Appendix~\ref{subsec:implementation-detail} for
invocation examples and metadata format.

\paragraph{Skill Extensibility.}
The \texttt{skill-maker} meta-skill enables creation of new domain-specific
skills following established patterns. This allows SPOQ deployments to
accumulate organizational knowledge as reusable agent capabilities.

\newpage
\subsection{Integration: From Specification to Execution}
\label{subsec:integration}

The implementation components integrate as follows:

\begin{enumerate}
    \item \textbf{Epic Creation:} Human uses \texttt{/epic-planning} to decompose
    a goal. Output: \texttt{EPIC.md} + task YAML files in \texttt{spoq/epics/active/}.

    \item \textbf{Validation:} \texttt{/epic-validation} and \texttt{/task-validation}
    score specifications against metrics. Failed validations trigger revision.

    \item \textbf{Execution:} \texttt{/agent-execution} reads task files, computes
    waves, and dispatches Opus workers. Each agent receives task YAML as context.

    \item \textbf{Tracking:} Agents write journal entries on session completion.
    The journal accumulates work history for the epic.

    \item \textbf{Code Validation:} \texttt{/agent-validation} scores completed
    tasks. Failed tasks queue for rework with remediation guidance.

    \item \textbf{Completion:} On all tasks passing, the epic moves to
    \texttt{spoq/epics/complete/}.
\end{enumerate}

This lifecycle ensures traceability from initial goal through validated
delivery, with quality gates preventing propagation of defects.
Appendix~\ref{subsec:integration-detail} discusses integration patterns
for CI/CD pipelines, project management tools, and git workflows.

\paragraph{Repository Bootstrap.}
To reduce adoption friction, SPOQ provides cross-platform installer
scripts (\texttt{spoq-init.sh} for Linux and macOS, \texttt{spoq-init.ps1}
for Windows) that automate repository initialization. In fresh-install mode,
the scripts create the five-category directory structure
(\texttt{code/}, \texttt{documents/}, \texttt{spoq/}, \texttt{infrastructure/},
\texttt{tests/}), copy the full skill definitions described in
Section~\ref{subsec:skills}, and generate starter \texttt{CLAUDE.md} and
\texttt{journal.md} files. An optional \texttt{-{}-full} flag provisions an
example epic so that new adopters can exercise the pipeline immediately.

For repositories that already use an earlier directory layout, the
\texttt{-{}-upgrade} flag activates a migration path. The upgrade routine
detects legacy structures (e.g., \texttt{automation/tasks/}), relocates epics
to \texttt{spoq/epics/active/} with conflict-aware merging, refreshes
skill definitions from the canonical source while preserving timestamped
backups, and updates configuration references. Together, these modes allow a
team to adopt or migrate to SPOQ with a single command invocation rather than
manual directory scaffolding.
